%% kapa: exam=1 round=33 year=2022 period=1 form=A subject=real_estate subject_ko=부동산학원론 seq=21 number=101
%% source: Q-Net 감정평가사 기출문제 (cst00309) artlSeq=5209997 file=raw/1008/5209997/2022년도 제33회 감정평가사 1차시험 1교시(민법,경제학,부동산학).pdf
%% conversion: kordoc 4.24.51 -> tools/gen1.py (auto); review=proofread-sonnet; flags=none
\begin{question}{101}
부동산금융에 관한 설명으로 옳지 않은 것은? (단, 주어진 조건에 한함)
\begin{choices}
  \item 대출채권의 듀레이션(평균 회수기간)은 만기일시상환대출이 원리금균등분할상환대출보다 길다.
  \item 대출수수료와 조기상환수수료를 부담하는 경우 차입자의 실효이자율은 조기상환시점이 앞당겨질수록 상승한다.
  \item 금리하락기에 변동금리대출은 고정금리대출에 비해 대출자의 조기상환위험이 낮다.
  \item 금리상승기에 변동금리대출의 금리조정주기가 짧을수록 대출자의 금리위험은 낮아진다.
  \item 총부채원리금상환비율(DSR)과 담보인정비율(LTV)은 소득기준으로 채무불이행위험을 측정하는 지표이다.
\end{choices}
\end{question}
\answer{5}
